\documentclass[11pt]{article}
\usepackage{amsmath,amssymb,amsthm}
\usepackage{booktabs}
\usepackage[hidelinks]{hyperref}
\emergencystretch=1.5em

\title{Direct Answer and Evidence: GREEN-028\\
\large Gowers Box Norms over Finite Fields}
\author{}
\date{}

\begin{document}
\maketitle

\section*{Direct Answer}
\noindent\fbox{\parbox{0.97\textwidth}{%
\textbf{Verdict: UNRESOLVED.} The requested correlation statement remains open; no proof and no counterexample are known.
}}\medskip

\noindent\textbf{Proof status:} not obtained.

\noindent\textbf{Counterexample status:} not found (numerical search for $p=5,7$ found no counterexample but is not exhaustive).

\noindent\textbf{Exact missing implication:} The Fourier-mass hypothesis on $B_f$ is not known to imply algebraic correlation with the infinite-dimensional target class; no known implication from the spectral condition to correlation exists.

\noindent\textbf{Scope:} General theorem for all odd primes $p$ and dimensions $n$; partial results hold where stated.

All numerical residuals quoted here were independently confirmed by an
external deterministic verification suite, which was independently revalidated at build time during the
preparation of this document.

\section*{Unconditional results}
\begin{description}
\item[R1.] \textbf{Spectral reformulation (Proposition A).} With $B_f(x_1,x_2,y_1,y_2)=f(x_1,y_1)\overline{f(x_2,y_1)}\,\overline{f(x_1,y_2)}f(x_2,y_2)$ on $G^4$ and $\Xi_0=\{(\xi_1,\xi_2,\xi_3,\xi_4)\in\widehat{G}^{\,4}:\ \xi_1+\xi_2+\xi_3+\xi_4=0\}$, one has $$\mathbb{E}_h\,\|\Delta_{(h,h)}f\|_\square^4\;=\;\sum_{\xi\in\Xi_0}\bigl|\widehat{B_f}(\xi)\bigr|^2.$$ The hypothesis says exactly that the four-point box product $B_f$ carries Fourier $L^2$-mass at least $\delta$ on the zero-sum-frequency subsystem $\Xi_0$, of codimension $n$ in $\widehat{G^4}$. Machine-checked at $(p,n)=(3,1),(3,2),(5,1)$ with residuals between $1.11\times10^{-16}$ and $6.11\times10^{-16}$ and Parseval diagnostic total Fourier mass equal to $1$ at every size.
\item[R2.] \textbf{Multiplicative calculus of the target class.} (i) For $n=1$ every quadratic phase absorbs into the three one-variable factors: $Q(x,y)=\alpha(x)+\beta(y)+\gamma(x+y)+\kappa$ via $xy=\bigl((x+y)^2-x^2-y^2\bigr)/2$, valid for odd $p$. (ii) For $n\ge2$ the monomial $x_1y_2$ admits no such representation (mixed finite differences force $w_1s_2=s_1w_2$ for all $s,w$, false at $s=e_2$, $w=e_1$), so the phase genuinely enlarges the class. (iii) Factorization: $\|a(x)b(y)w(x+y)\|_\square^4=\|w\|_{U^2(G)}^4$ for unimodular $a,b$. (iv) Reduction: for $f=a(x)b(y)c(x+y)$, $$\mathbb{E}_h\|\Delta_{(h,h)}f\|_\square^4=\mathbb{E}_h\|\Delta_{2h}c\|_{U^2(G)}^4=\|c\|_{U^3(G)}^8,$$ so inside the class the hypothesis is precisely a Gowers $U^3$ condition, and the inverse theorem for $U^3$ over $\mathbb{F}_p^n$ explains every class member obeying it by a quadratic phase. The reduction identity was measured to residual $2.220\times10^{-16}$ at $(p,n)=(3,1)$.
\item[R3.] \textbf{Diagonal family and twisted trilinear inequality (Observation C).} For every unimodular $d\colon G\to\mathbb{C}$, $f(x,y)=d(x-y)$ satisfies $\Delta_{(h,h)}f\equiv 1$, hence meets the hypothesis with $\delta=1$. Under the bijection $x=(m+s)/2$, $y=(m-s)/2$ (valid since $p$ is odd), correlation against the target class becomes $$\Bigl|\;\mathbb{E}_{s,m}\;d(s)\,\overline{a\bigl(\tfrac{m+s}{2}\bigr)}\,\overline{b\bigl(\tfrac{m-s}{2}\bigr)}\,\overline{c(m)}\;e_p\!\bigl(-Q\bigl(\tfrac{m+s}{2},\tfrac{m-s}{2}\bigr)\bigr)\;\Bigr|;$$ at $n=1$ the twist absorbs by (ii), so the maximal-hypothesis case of the problem is exactly a uniform positive lower bound for this trilinear correlation over all unimodular $d$, which is open.
\item[R4.] \textbf{Rigidity: sufficiency half of rank-one diagonal increments (Theorem D).} If $f=a(x)b(y)d(x-y)e_p(Q)$ with $Q(z)=\tfrac12 z^{T}Mz+l^{T}z+\kappa$ quadratic on $G\times G$, then $$\Delta_{(h,h)}f(x,y)=\alpha_h(x)\,\beta_h(y)\,e_p(\omega_h),$$ with $\alpha_h(x)=a(x+h)\overline{a(x)}\,e_p\bigl(((M_{xx}{+}M_{xy})h)^{T}x\bigr)$, $\beta_h(y)=b(y+h)\overline{b(y)}\,e_p\bigl(((M_{yx}{+}M_{yy})h)^{T}y\bigr)$, $\omega_h=\tfrac12(h,h)^{T}M(h,h)+l^{T}(h,h)$; verified to maximal deviation $5.118\times10^{-16}$ over all $(h,x,y)$ at $(p,n)=(5,1)$. Sharpness: multiplication by the cubic mixed phase $e_p(xy^2)$ destroys rank-one factorization for $4$ of $5$ directions (fraction $0.8$). The converse classification reduces to the multiplicative parallelogram equation $\Theta(u+w,v+t)\Theta(u-w,v-t)=\Theta(u+t,v-w)\Theta(u-t,v+w)$ and remains open. Structurally, the natural rigidity family $\{a\,b\,d(x-y)e_p(Q)\}$ differs from the transmitted class $\{a\,b\,c(x+y)(-1)^q\}$ once $n\ge2$, since $e_p(\lambda x_1y_2)$ belongs to the former but not to the absorbable span of the latter.
\item[R5.] \textbf{Small-field numerics: degrees-of-freedom-dominated regime.} Maximizing $M(d)=\sup\bigl|\mathbb{E}_{x,y}\,d(x-y)\overline{a(x)}\,\overline{b(y)}\,\overline{c(x+y)}\bigr|$ over unimodular factors (which dominates the full stated class when $n=1$) by alternating phase maximization from $1000$ random restarts: for $d=e_p(s^3)$ the best correlation is $0.7236067977499792$ at $p=5$ (mean $0.7225012249409789$, sd $0.01745$) and $0.6772769730217724$ at $p=7$ (mean $0.6758013999033685$, sd $0.01399$); random-exponent controls give $0.7236067977499792$ ($p=5$) and $0.8160210786702251$ ($p=7$); the control $d=e_p(2s^2)$ attains correlation exactly $1$ via the witness $e_p(4x^2)e_p(4y^2)e_p(-2s^2)$ (a naive witness yields the Gauss-sum value $1/\sqrt{p}$, recorded as erratum). At $n=1$, $p\in\{5,7\}$, optima cluster in $0.68$--$0.82$ irrespective of structure: $3p$ unimodular degrees of freedom dominate $p^2$ constraints, so these sizes exhibit no separation between structured and random diagonal data and give neither evidence for nor against asymptotic decay.
\end{description}

\section*{Derivation}
% LaTeX-ready body text for the Derivation section.
% Contract: this file contains pure mathematical LaTeX (no code, no logs, no
% paths); it is embedded verbatim.

\subsection{Setup and conventions}
Throughout, $p$ is an odd prime, $G=\mathbb{F}_p^n$, $e_p(t)=\exp(2\pi i t/p)$,
and $\mathbb{E}$ denotes the uniform average. The diagonal derivative and box
norm are
\[
\Delta_{(h,h)}f(x,y)=f(x+h,y+h)\overline{f(x,y)},\qquad
\|g\|_\square^4=\mathbb{E}_{x_1,x_2,y_1,y_2}\,
g(x_1,y_1)\overline{g(x_2,y_1)}\,\overline{g(x_1,y_2)}\,g(x_2,y_2).
\]
The target class under study is
$\{\,a(x)b(y)c(x+y)e_p(Q(x,y)) : a,b,c \text{ unimodular},\ Q \text{ quadratic}\,\}$,
which contains both readings of $(-1)^{q(x,y)}$ as subclasses; every result
below addresses this larger class.

\subsection{Step 1: the hypothesis is a spectral mass condition}
Let $B_f(x_1,x_2,y_1,y_2)=f(x_1,y_1)\overline{f(x_2,y_1)}\,
\overline{f(x_1,y_2)}f(x_2,y_2)$. Expanding the box expression of
$\Delta_{(h,h)}f$ term by term, the four shifted factors assemble to
$B_f(z+(h,h,h,h))$ and the four conjugated unshifted factors to
$\overline{B_f(z)}$, so with $u=(h,h,h,h)$,
\[
\|\Delta_{(h,h)}f\|_\square^4=\mathbb{E}_{z\in G^4}B_f(z+u)\overline{B_f(z)}
=\sum_{\xi\in\widehat{G^4}}\bigl|\widehat{B_f}(\xi)\bigr|^2 e_p(\xi\cdot u).
\]
Averaging over $h$ retains exactly the characters with
$\xi_1+\xi_2+\xi_3+\xi_4=0$:
\begin{equation}\label{eq:spectral}
\mathbb{E}_h\,\|\Delta_{(h,h)}f\|_\square^4
=\sum_{\xi\in\Xi_0}\bigl|\widehat{B_f}(\xi)\bigr|^2 .
\end{equation}
The left side is nonnegative, and \eqref{eq:spectral} exhibits it as Fourier
$L^2$-mass of $B_f$ on the zero-sum subsystem $\Xi_0$, of codimension $n$ in
$\widehat{G^4}$.

\subsection{Step 2: multiplicative calculus of the class}
For odd $p$, $xy=\bigl((x+y)^2-x^2-y^2\bigr)/2$, whence for $n=1$ every
quadratic $Q(x,y)=Ax^2+Bxy+Cy^2+Dx+Ey+F$ splits as
$\alpha(x)+\beta(y)+\gamma(x+y)+\kappa$ with
$\alpha=(A-\tfrac{B}{2})x^2+Dx$, $\beta=(C-\tfrac{B}{2})y^2+Ey$,
$\gamma=\tfrac{B}{2}s^2$, $\kappa=F$: at $n=1$ the phase is redundant. For
$n\ge2$ the monomial $x_1y_2$ obstructs: a mixed finite difference applied to
a purported identity $x_1y_2=\alpha+\beta+\gamma+\kappa$ yields
$\gamma(s+w)-\gamma(s)=w_1s_2$ up to constants, and symmetrizing in $(s,w)$
forces $w_1s_2=s_1w_2$ for all $s,w$ — false at $s=e_2,w=e_1$. Next, against
the box pattern the factors $a(x)b(y)$ cancel pairwise:
\[
\|a(x)b(y)w(x+y)\|_\square^4=\|w\|_{U^2(G)}^4 ,
\]
because $s_{ij}=x_i+y_j$ sweep every additive quadruple exactly once.
Consequently, for $f=a(x)b(y)c(x+y)$,
\[
\mathbb{E}_h\|\Delta_{(h,h)}f\|_\square^4
=\mathbb{E}_h\|\Delta_{2h}c\|_{U^2(G)}^4=\|c\|_{U^3(G)}^8 :
\]
inside the target class the hypothesis is precisely a Gowers $U^3$ condition,
and the inverse theorem for $U^3$ over $\mathbb{F}_p^n$ explains every class
member obeying it by a quadratic phase — confirming degree $2$ as the right
target complexity within the class.

\subsection{Step 3: the diagonal family and an equivalent inequality}
Every $f=d(x-y)$ satisfies $\Delta_{(h,h)}f=d(x-y)\overline{d(x-y)}=1$, so the
hypothesis holds with $\delta=1$. Under the bijection
$x=(m+s)/2$, $y=(m-s)/2$, correlation of such $f$ with the class becomes
\[
\Bigl|\;\mathbb{E}_{s,m}\;
d(s)\,\overline{a\bigl(\tfrac{m+s}{2}\bigr)}\,
\overline{b\bigl(\tfrac{m-s}{2}\bigr)}\,\overline{c(m)}\,
e_p\!\bigl(-Q\bigl(\tfrac{m+s}{2},\tfrac{m-s}{2}\bigr)\bigr)\;\Bigr| ;
\]
at $n=1$ the twist absorbs by Step 2, so the maximal-hypothesis case of the
problem is exactly the open problem of a uniform positive lower bound for this
trilinear correlation over all unimodular $d$.

\subsection{Step 4: rigidity of rank-one diagonal increments}
If $f=a(x)b(y)d(x-y)e_p(Q)$ with $Q(z)=\tfrac12 z^{T}Mz+l^{T}z+\kappa$ on
$G\times G$, then $Q(z+\delta)-Q(z)=(M\delta)^{T}z+\tfrac12\delta^{T}M\delta
+l^{T}\delta$ has vanishing mixed $xy$ coefficient, while
$d(x+h-(y+h))\overline{d(x-y)}=1$; hence
\[
\Delta_{(h,h)}f(x,y)
=a(x+h)\overline{a(x)}\,e_p\bigl((Ph)^{T}x\bigr)\cdot
 b(y+h)\overline{b(y)}\,e_p\bigl((Rh)^{T}y\bigr)\cdot e_p(\omega_h),
\]
with $P=M_{xx}+M_{xy}$, $R=M_{yx}+M_{yy}$,
$\omega_h=\tfrac12(h,h)^TM(h,h)+l^T(h,h)$: pointwise rank-one increments.
Multiplication by $e_p(xy^2)$ contributes the mixed term $2hxy$ and destroys
the splitting for every $h\ne0$. The converse reduces to the parallelogram
equation
$\Theta(u+w,v+t)\Theta(u-w,v-t)=\Theta(u+t,v-w)\Theta(u-t,v+w)$ and remains
open; the natural rigidity class $\{a\,b\,d(x-y)e_p(Q)\}$ differs from the
transmitted class $\{a\,b\,c(x+y)(-1)^q\}$ once $n\ge2$, since
$e_p(\lambda x_1y_2)$ belongs to the former but not to the absorbable span of
the latter.

\subsection{Step 5: what is missing}
The spectral condition of Step 1 controls only aggregate Fourier mass of
$B_f$ on $\Xi_0$, whereas the conclusion demands correlation with members of
an infinite-dimensional threefold-factor class. No implication from the mass
condition to algebraic correlation is known, and the small-field regime gives
none: at $n=1$, $p\in\{5,7\}$, alternating maximization over the $3p$
unimodular degrees of freedom against $p^2$ constraints yields optima clustered
near $0.68$--$0.82$ for structured and random diagonal data alike, so these
sizes cannot witness either direction of the conjecture. Closing the gap
requires either a structure theorem extracting class membership from
$\Xi_0$-mass (one route: settle the parallelogram-equation converse), or a
counterexample family with decaying optimal correlation (related to slice and
partition rank of the kernel $(s,m)\mapsto d(s)e_p(Q(s,m))$).


\section*{Evidence}
% LaTeX-ready body text for the Evidence section.
% Contract: this file contains pure mathematical LaTeX (no code, no logs, no
% paths); it is embedded verbatim.

All quantitative claims below were obtained by deterministic computation in
exact rational, modular, and high-precision floating-point arithmetic; every
identity was recomputed from its definition on both sides rather than asserted.

\subsection{Identity checks}
\begin{itemize}
\item \textbf{Spectral reformulation.} With $\{\pm1\}$-valued random data the
two sides of \eqref{eq:spectral} were computed independently: at $(p,n)=(3,1)$
the values $0.6049382716049383$ and $0.6049382716049382$ agree to
$1.11\times10^{-16}$; at $(3,2)$, $0.307422648986435$ and
$0.3074226489864356$ agree to $6.11\times10^{-16}$; at $(5,1)$, $0.45728$ and
$0.4572800000000002$ agree to $1.67\times10^{-16}$. The Parseval diagnostic
$\sum_{\xi\in\widehat{G^4}}|\widehat{B_f}(\xi)|^2=1$ held at all three sizes.
\item \textbf{Class reduction.} At $(p,n)=(3,1)$ with random unimodular
factors, both sides of
$\mathbb{E}_h\|\Delta_{(h,h)}[abc]\|_\square^4=\|c\|_{U^3(G)}^8$ evaluated to
$0.9999999999999997$ and $0.9999999999999999$, agreeing to
$2.220\times10^{-16}$.
\item \textbf{Absorption.} Gaussian elimination over $\mathbb{F}_p$ produced
consistent coefficient systems with zero pointwise mismatches on random tests
at $(p,n)=(3,1),(5,1),(5,2)$ for symmetric cross-blocks; a forced asymmetric
cross-block at $(5,2)$ produced an inconsistent system whose rank falls one
equation short, confirming the obstruction of Step 2 quantitatively.
\item \textbf{Rigidity splitting.} At $(p,n)=(5,1)$ the explicit rank-one
splitting held over all triples $(h,x,y)$ with maximal deviation
$5.118\times10^{-16}$ from the product of its two factors; after multiplication
by $e_p(xy^{2})$, four of the five shifts admitted no rank-one factorization
(fraction $0.8$, measured by the $2\times2$ minor test).
\end{itemize}

\subsection{Correlation experiment}
The maximal-hypothesis correlation
$M(d)=\sup\bigl|\mathbb{E}_{x,y}\,d(x-y)\overline{a(x)}\,\overline{b(y)}\,
\overline{c(x+y)}\bigr|$ over unimodular factors was optimized by alternating
phase maximization from one thousand random restarts per configuration:
\begin{center}
\begin{tabular}{llr}
\toprule
target $d$ & statistic & value \\
\midrule
$e_5(s^{3})$ & best / mean / sd & $0.7236067977499792$ / $0.7225012249409789$ / $0.01745$ \\
$e_7(s^{3})$ & best / mean / sd & $0.6772769730217724$ / $0.6758013999033685$ / $0.01399$ \\
random exponents, $p=5$ & best & $0.7236067977499792$ \\
random exponents, $p=7$ & best & $0.8160210786702251$ \\
$e_p(2s^{2})$ (control) & exact correlation & $1$ exactly \\
\bottomrule
\end{tabular}
\end{center}
The quadratic control attains correlation exactly $1$ through the witness
$a=e_p(4x^{2})$, $b=e_p(4y^{2})$, $c=e_p(-2s^{2})$, using
$4x^{2}+4y^{2}-2(x+y)^{2}=2(x-y)^{2}$; a commonly guessed naive witness yields
instead the Gauss-sum value $1/\sqrt{p}$, an instructive erratum recorded here.

Interpretation: structured cubic targets behave indistinguishably from random
targets at these sizes; the regime is dominated by the count of free
parameters rather than by arithmetic structure, so the experiment provides no
counterexample and no separation — consistent with, but not evidential for,
either direction of the open problem.


\begin{thebibliography}{1}
\bibitem{unsolvedmath} UnsolvedMath problem archive, entry \texttt{GREEN-028},
\url{https://www.unsolvedmath.com/problems/GREEN-028}.
\end{thebibliography}

\end{document}
